\chapter{Indradhanu Heavy: capacity, a fair control and the haze}\label{ch:heavy}

\begin{onething}
Tripling the encoder's width multiplies its parameters by about nine, because almost every weight sits in a
$C\times C$ matrix; the resulting 33.9\,M-parameter Indradhanu Heavy, re-pretrained and fine-tuned beside a
control given exactly the same recipe, beat the control and the shipped model on detection at every
neighbouring checkpoint, once a thin-cloud mask that fired on haze was left out, while it reconstructed
no better than the control.
\end{onething}

This chapter closes Part VIII with the experiment that finally moved detection. It has four strands: the
arithmetic of width (why parameters grow like width squared, counted to the last parameter); the pretraining run
on a rented H100 and the two operational traps it hit; the logic of a fair comparison (a control, and three
neighbouring checkpoints); and a mask head that called haze ``cloud'' and so deleted most of a scene's targets.
Its lesson is the one \cref{ch:uncertainty,ch:decoderwork} pointed towards: judge detectors on detection.

Provenance tags abbreviate paths:
\begin{itemize}
\item \file{F002-heavy/} is under \file{research/experiments/indradhanu-se-pretrain/};
\item \file{FTH1-heavy-joint/} and \file{FTH2-heavy-joint-more/},
\item \file{FTH3-heavy-joint-more2/} and \file{FTB1-base-joint/}:
  all four are under\\ \file{research/experiments/indradhanu-finetune/}.
\end{itemize}

\begin{analogy}
You want to know whether a new engine makes a car faster. Last year's car won with a driver who practised the
track for a month. Put the new engine in a car, hand it to a driver who has had one afternoon, and if it loses,
you have learned nothing about the engine. The fair test is two cars, old engine and new, with the same driver,
the same afternoon of practice, the same tyres. Here the engine is the pretrained model (Heavy against F001),
the month of practice is FT04's four-stage curriculum, and the afternoon is a single joint fine-tuning run that
both engines get.
\end{analogy}

\section{Why capacity was the untested question}\label{sec:heavy:why}

Look back at \cref{tab:decoderwork:cumulative}. Cloud exclusion was a real gain; FT07 and FT08 cost a little;
FT09b gained by a hair and FT09c showed the gain was a lucky checkpoint. Every one of those runs \emph{kept
FT04's encoder at the same size}: widths 32, 64, 160 and 256 across its four stages. \Cref{ch:decoderwork} found
the limit they could not move: each stage-1 token squeezes 512 numbers into 32. FT09 attacked the decoder side of
that; it did not change how much the encoder can hold. As the F002 README puts it: ``A wider encoder needs
re-pretraining'' \tagM{F002-heavy/README.md}. You cannot widen a trained encoder and keep its weights, because
every weight matrix changes shape.

\begin{definition}[Capacity and width]
The \emph{capacity}\index{capacity} of a model is how much it can represent. For a network of fixed depth it
grows with its \emph{width}\index{width (of a layer)}: the number of channels each layer carries.
\end{definition}

\section{The architecture}\label{sec:heavy:arch}

\begin{table}[htbp]
\centering\small
\begin{tabular}{llll}
\toprule
setting & F001 / FT04 & Heavy & ratio\\
\midrule
\texttt{embed\_dims} (stage widths) & 32, 64, 160, 256 & 96, 192, 480, 768 & 3$\times$\\
\texttt{num\_heads} & 2, 2, 5, 8 & 3, 6, 10, 12 & \\
channels per head & 16, 32, 32, 32 & 32, 32, 48, 64 & \\
spectral latent (\texttt{compressed\_channels}) & 32 & 64 & 2$\times$\\
\texttt{token\_dim} & 32 & 32 & 1$\times$\\
decoder & SegFormer, dim 256 & progressive, dim 512 & \\
decoder stage dims & -- & 256 / 128 & \\
reduction ratios; blocks & 8 4 2 1;\ 2 2 2 2 & same & \\
mask ratio & 0.65 & 0.65 & \\
\bottomrule
\end{tabular}
\caption{F001 (FT04's pretrained parent, \cref{ch:pretrain}) against Heavy, both built by
\texttt{SpectralStemSegFormerMAE}. Source: \tagM{F002-heavy/config.json}.}
\label{tab:heavy:arch}
\end{table}

\begin{sourcenote}[: ``4x wide'' is a name, the widths are 3x]
The run was launched on the user's ``Go, but 4x wide''. The course (part 68) records that the README title, the
config's \texttt{\_description} and the commit all said ``4x wide''. The configuration has encoder widths
$96/32=3$ and $768/256=3$ times FT04's, and the spectral latent and decoder 2 times; nothing in the model is 4
times. The source now agrees: the README title reads ``3x-wide pretraining'', with ``implemented as 3x encoder
widths \ldots{} which this README first called 4x (corrected 2026-10-29)'', and the \texttt{\_description} reads
``3x-wide (encoder)'' (\file{research/experiments/indradhanu-se-pretrain/F002-heavy/README.md},
\file{config.json}). Read ``4x'' as the launch instruction, not a measurement.
\end{sourcenote}

The decoder is FT09's progressive decoder (\cref{ch:decoderwork}). It did not help FT04 on its own; here it is
pretrained from the start, so it never has the warm-start problem. \texttt{token\_dim} stays 32 because ``the
reduced stem's wavelength encoding is 32 wide''.

\subsection{The stage-1 squeeze, three times wider}

\begin{align*}
\text{FT04} &: 16\text{ pixels}\times32\text{ channels} = 512 \text{ in};\ 32\text{ out};\ 512/32 = 16\text{ to }1;\ 32/16 = 2\text{ per pixel}\\
\text{Heavy} &: 16\times64 = 1{,}024 \text{ in};\ 96\text{ out};\ 1{,}024/96 = 10.7\text{ to }1;\ 96/16 = 6\text{ per pixel}
\end{align*}

\section{Why parameters grow like width squared}\label{sec:heavy:width}

\subsection{One encoder block, as a formula}

A transformer block (\cref{ch:block}) is LayerNorm, efficient self-attention with spatial reduction ratio $R$
(\cref{ch:esa}), LayerNorm, Mix-FFN with expansion 4. Count its parameters for width $C$, layer by layer. A
linear layer from $a$ to $b$ channels has $ab+b$ parameters; a LayerNorm over $C$ channels has $2C$ (a weight and a
bias each).
\begin{align*}
\text{norm1} &: 2C\\
q,k,v,\text{proj} &: 4\,(C\cdot C + C) = 4C^2 + 4C\\
\text{reduction conv } (R\times R,\ \text{if } R>1) &: C\cdot C\cdot R^2 + C = R^2C^2 + C\\
\text{reduction norm (if } R>1) &: 2C\\
\text{norm2} &: 2C\\
\text{fc1 } (C\to4C) &: C\cdot4C + 4C = 4C^2 + 4C\\
\text{depthwise }3\times3\text{ conv on }4C &: 4C\cdot1\cdot9 + 4C = 40C\\
\text{fc2 } (4C\to C) &: 4C\cdot C + C = 4C^2 + C
\end{align*}
Collect the $C^2$ terms: attention $4C^2$, reduction $R^2C^2$, FFN $4C^2+4C^2=8C^2$. Collect the $C$ terms:
$2+4+1+2+2+4+40+1=56$. So
\begin{equation}\label{eq:heavy:block}
P_\text{block}(C,R) = (12+R^2)\,C^2 + 56\,C \quad (R>1),\qquad P_\text{block}(C,1) = 12\,C^2 + 53\,C,
\end{equation}
where for $R=1$ the reduction convolution and its norm are absent ($56-1-2=53$).

\begin{worked}[label={we:heavy:block1}]{Stage-1 block, FT04 and Heavy}
Stage 1 has $R=8$, so the $C^2$ coefficient is $12+64=76$. Do it on paper.
\begin{align*}
\text{FT04, } C=32 &: 76\times32^2 + 56\times32 = 76\times1{,}024 + 1{,}792 = 77{,}824 + 1{,}792 = 79{,}616\\
\text{Heavy, } C=96 &: 76\times96^2 + 56\times96 = 76\times9{,}216 + 5{,}376 = 700{,}416 + 5{,}376 = 705{,}792\\
\text{ratio} &: 705{,}792/79{,}616 = 8.86
\end{align*}
Line by line for Heavy: norm1 192; q,k,v,proj $4\times(9{,}216+96)=37{,}248$; reduction
$96\times96\times64+96=589{,}920$; reduction norm 192; norm2 192; fc1 $96\times384+384=37{,}248$; dwconv
$384\times9+384=3{,}840$; fc2 $384\times96+96=36{,}960$. Sum: $192+37{,}248+589{,}920+192+192+37{,}248+3{,}840+36{,}960=705{,}792$.
The $8\times8$ reduction convolution alone is $589{,}920/705{,}792=84\%$ of the block.
\end{worked}

\paragraph{Why the ratio is a little under 9.} Tripling $C$ multiplies the $C^2$ term by $3^2=9$ and the $C$
term by only 3:
\begin{align*}
\frac{P(3C)}{P(C)} = \frac{9\,aC^2 + 3\cdot56\,C}{aC^2 + 56\,C}
 = 9 - \frac{6\cdot56\,C}{aC^2+56\,C}.
\intertext{For stage 1 of FT04 ($a=76$, $C=32$):}
\frac{6\times56\times32}{79{,}616} = \frac{10{,}752}{79{,}616} = 0.135,\qquad 9-0.135 = 8.865 .
\end{align*}
The wider the layer, the smaller the correction: width scaling is ``about width squared'' because the $C^2$
terms dominate.

\begin{table}[htbp]
\centering\small
\begin{tabular}{lrrrr}
\toprule
stage ($C_\text{FT04}\to C_\text{Heavy}$, $R$) & FT04 block & Heavy block & ratio & formula check (Heavy)\\
\midrule
1 ($32\to96$, 8) & 79,616 & 705,792 & 8.86 & $76\cdot9{,}216+5{,}376$\\
2 ($64\to192$, 4) & 118,272 & 1,042,944 & 8.82 & $28\cdot36{,}864+10{,}752$\\
3 ($160\to480$, 2) & 418,560 & 3,713,280 & 8.87 & $16\cdot230{,}400+26{,}880$\\
4 ($256\to768$, 1) & 800,000 & 7,118,592 & 8.90 & $12\cdot589{,}824+40{,}704$\\
\bottomrule
\end{tabular}
\caption{One block per stage. Counts from building both models on CPU \tagC
(\texttt{tools/compute\_39-41\_output.txt}); every one matches \cref{eq:heavy:block}.}
\label{tab:heavy:blocks}
\end{table}

\begin{figure}[htbp]
\centering
\begin{tikzpicture}
\begin{axis}[width=0.75\linewidth,height=5.5cm,xlabel={width multiplier $k$},ylabel={$P(kC)/P(C)$},
  xmin=1,xmax=4,ymin=0,ymax=16,domain=1:4,samples=60,grid=major,grid style={gray!20},
  legend pos=north west,legend cell align=left]
\addplot[loss,thick] {(76*1024*x^2+56*32*x)/(76*1024+56*32)};
\addlegendentry{stage-1 block, block formula}
\addplot[maskgrey,dashed] {x^2};
\addlegendentry{$k^2$}
\addplot[only marks,mark=*,attn] coordinates {(3,8.865)};
\node[attn,anchor=west,font=\small] at (axis cs:3.05,7.6) {Heavy: 8.86};
\end{axis}
\end{tikzpicture}
\caption{Parameter growth of a stage-1 block with width, drawn to scale. The curve hugs $k^2$; the gap is the
linear (bias and norm) terms.}
\label{fig:heavy:growth}
\end{figure}

\subsection{The whole encoder, counted}

The patch embeddings (\cref{ch:patchembed}): stage 1 is a $4\times4$ convolution with stride 4 from the spectral
latent; stages 2--4 are $3\times3$, stride 2, from the previous stage. Each is followed by a LayerNorm.
\begin{align*}
\text{stage 1} &: 64\times96\times16 + 96 + 2\times96 = 98{,}304 + 96 + 192 = 98{,}592\\
\text{stage 2} &: 96\times192\times9 + 192 + 384 = 165{,}888 + 576 = 166{,}464\\
\text{stage 3} &: 192\times480\times9 + 480 + 960 = 829{,}440 + 1{,}440 = 830{,}880\\
\text{stage 4} &: 480\times768\times9 + 768 + 1{,}536 = 3{,}317{,}760 + 2{,}304 = 3{,}320{,}064\\
\text{sum} &: 98{,}592+166{,}464 = 265{,}056;\ +830{,}880 = 1{,}095{,}936;\ +3{,}320{,}064 = 4{,}416{,}000
\end{align*}
(FT04's stage-1 embedding: $32\times32\times16+32+64=16{,}480$.) Two blocks per stage:
\begin{align*}
2\times705{,}792 &= 1{,}411{,}584\\
2\times1{,}042{,}944 &= 2{,}085{,}888\\
2\times3{,}713{,}280 &= 7{,}426{,}560\\
2\times7{,}118{,}592 &= 14{,}237{,}184\\
\text{blocks} &: 1{,}411{,}584+2{,}085{,}888 = 3{,}497{,}472;\ +7{,}426{,}560 = 10{,}924{,}032\\
&\phantom{:}\ 10{,}924{,}032+14{,}237{,}184 = 25{,}161{,}216
\end{align*}
and a LayerNorm after each stage, $2\times(96+192+480+768)=2\times1{,}536=3{,}072$. The encoder:
\[
25{,}161{,}216 + 4{,}416{,}000 + 3{,}072 = 29{,}580{,}288 .
\]

\subsection{The decoder and the whole model}

\begin{worked}[label={we:heavy:decoder}]{Heavy's progressive decoder, 4,279,616 parameters}
Widths: stages $[96,192,480,768]$, \code{decoder_dim} 512, skip 64 channels (the spectral latent), stage dims
(256, 128), 64 output channels.
\begin{align*}
\text{projections} &: (96+192+480+768)\times512 + 4\times512 = 786{,}432 + 2{,}048 = 788{,}480\\
\text{fuse} &: 2{,}048\times512 + 512 = 1{,}048{,}576 + 512 = 1{,}049{,}088\\
\texttt{up\_half}\ 1 &: (512+64)\times256\times9 + 256 = 1{,}327{,}104 + 256 = 1{,}327{,}360\\
\texttt{up\_half}\ 2 &: 256\times256\times9 + 256 = 589{,}824 + 256 = 590{,}080\\
\texttt{up\_full}\ 1 &: (256+64)\times128\times9 + 128 = 368{,}640 + 128 = 368{,}768\\
\texttt{up\_full}\ 2 &: 128\times128\times9 + 128 = 147{,}456 + 128 = 147{,}584\\
\texttt{head} &: 128\times64 + 64 = 8{,}192 + 64 = 8{,}256\\[3pt]
\text{running total} &: 788{,}480+1{,}049{,}088 = 1{,}837{,}568;\ +1{,}327{,}360 = 3{,}164{,}928\\
&\phantom{:}\ +590{,}080 = 3{,}755{,}008;\ +368{,}768 = 4{,}123{,}776\\
&\phantom{:}\ +147{,}584 = 4{,}271{,}360;\ +8{,}256 = 4{,}279{,}616
\end{align*}
\end{worked}

\begin{table}[htbp]
\centering
\begin{tabular}{lrr}
\toprule
part & FT04 size (no seg head) & Heavy\\
\midrule
\texttt{spectral\_in} (stem) & 14,721 & 17,825\\
\texttt{norm} (BatchNorm, $2D$) & 64 & 128\\
encoder & 3,331,072 & 29,580,288\\
decoder & 2,164,736 & 4,279,616\\
\texttt{spectral\_out} & 4,257 & 6,337\\
\midrule
total & \textbf{5,514,850} & \textbf{33,884,194}\\
\bottomrule
\end{tabular}
\caption{Whole-model counts. The encoder and decoder rows are derived by hand above, Heavy's stem rows in
\cref{sec:heavy:stemcount}; all match the built models \tagC.}
\label{tab:heavy:total}
\end{table}

Add Heavy's column: $17{,}825+128=17{,}953$; $+29{,}580{,}288=29{,}598{,}241$; $+4{,}279{,}616=33{,}877{,}857$;
$+6{,}337=33{,}884{,}194$. That is the README's ``33.9\,M parameters (encoder 29.6\,M, decoder 4.3\,M) vs FT04
5.9\,M''. The 5.9\,M is the shipped FT04 \emph{with} its segmentation head (5,886,451, \cref{ch:shipping});
without the head, $33{,}884{,}194/5{,}514{,}850=6.1$. \textbf{Heavy is about six times FT04, with nine times its
encoder} ($29{,}580{,}288/3{,}331{,}072=8.9$).

\begin{sourcenote}[: every count verified]
Building both models on CPU (\file{research/textbook/tools/compute_39-41.py}) gives Heavy 33,884,194 total,
encoder 29,580,288, decoder 4,279,616, \code{encoder.patch_embeds[0]} 98,592, one stage-1 block 705,792; the
FT04-sized model 5,514,850, stage-1 block 79,616 \tagC.
\end{sourcenote}

%======================================================================
\section{One forward pass of Heavy, at full size}\label{sec:heavy:trace}

The counts above are a budget. This section follows the money: one batch through the model the product serves,
Indradhanu~V2 \texttt{v2.1.0} (FTH3 epoch~5, Heavy fine-tuned with its segmentation head, \cref{sec:heavy:product}),
at real sizes, with every tensor's shape. It is where \cref{ch:decoder}'s trace of the small model, \cref{ch:namer}'s
trace of the stem, and the counts above meet. The sizes:
\begin{itemize}
\item $B=2$ tiles of $128\times128$, the product's tile size (\code{default_patch_size} 128,
  \file{backend/allotrope/foundation_models/resolver.py});
\item the native bands of PRISMA bench scene set1 after \code{trained_band_mask} (below);
\item the training mask ratio $\rho=0.65$ (the same file), so $\lfloor1{,}024\times0.65\rfloor=665$ of the 1,024
  stage-1 tokens are hidden and 359 kept (\cref{ch:masking}).
\end{itemize}
Every shape in this section was recorded by forward hooks on the shipped checkpoint, loaded strictly into the real
class, by \file{research/textbook/tools/compute_heavy_trace.py}; its output is
\file{compute_heavy_trace_output.txt} \tagC. The script reads set1 exactly as the evaluator does
(\lstinline!load_prisma_scene(..., "native")!).

\subsection{How many bands: 185, all kept}

Set1's native cube after the product's band filter has 185 bands, 453.4 to 2,477.1\,nm \tagC. The model then drops
any band farther than 15\,nm from every wavelength it has trained on (\code{trained_band_mask},
\cref{ch:specdecoder}). Which wavelengths count as ``trained on'' is the buffer \code{trained_wl}, which starts as the
202-wavelength SpectralEarth curve and grows during fine-tuning.
\begin{center}\small
\begin{tabular}{@{}lrr@{}}
\toprule
record used & wavelengths in the record & set1 bands kept\\
\midrule
the pretraining curve alone (a pretrained model) & 202 & 165 (20 dropped)\\
v2.1.0's \code{trained_wl}, after fine-tuning & 733 & 185 (none dropped)\\
\bottomrule
\end{tabular}
\end{center}
The 20 bands a pretrained record drops are the three past 2,445\,nm (2,463.0 to 2,477.1), one at 1,958.6, seven
from 1,736.5 to 1,794.0, five from 1,459.3 to 1,502.0 and four from 1,306.2 to 1,339.1, the edges of
SpectralEarth's water-vapour gaps: $3+1+7+5+4=20$ \tagC. Fine-tuning on EnMAP (224 bands), PRISMA and AVIRIS-NG
(372 bands) filled those gaps in the record; it now runs to 2,500.4\,nm with no gap wider than 30\,nm.

\begin{sourcenote}[: 185 to 165 is a pretrained model's number]
\file{docs/lld/spectral-stem.md} (``set1 native: 185 -> 165'') and \cref{ch:specdecoder} quote 20 of 185 bands
dropped. That is right for a model whose record is the pretraining curve. The served v2.1.0 keeps all 185, because its
\code{trained_wl} has 733 entries \tagC. The FTH3 README's CPU timing (``1186 x 1196 x 185 bands'') is consistent with
the served model. So $N=185$ below.
\end{sourcenote}

\subsection{The stem, at 32,768 pixels}

The stem treats every pixel separately, so it flattens the batch to $P=B\times H\times W=2\times128\times128=32{,}768$
rows (\cref{ch:namer}, next, walks the same code step by step at the namer's sizes). The 185 kept bands fall into 45 of the 53 bins of
40\,nm, at most 6 to a bin \tagC.

\begin{table}[htbp]
\centering\small
\begin{tabular}{@{}lp{0.43\linewidth}l@{}}
\toprule
\# & operation & shape\\
\midrule
0 & reflectance in; normalised by the wavelength curve & \shape{(2, 185, 128, 128)}\\
1 & wavelength encoding of every band & \shape{(185, 32)}\\
2 & bin membership (one-hot) & \shape{(185, 53)}\\
3 & width-weighted values; weight per bin & \shape{(32768, 185)}; \shape{(32768, 53)}\\
4 & band vectors $\vect{k}(\lambda)$; their matrix $\mat{K}$; bin tokens & \shape{(185, 32)}; \shape{(185, 1696)};\\
 & & \shape{(32768, 53, 32)}\\
4 & bin position bias & \shape{(53, 32)}\\
5 & queries (64 of them, 4 heads of 8); keys and values & \shape{(32768, 4, 64, 8)}; \shape{(32768, 4, 53, 8)}\\
5 & attention weights (never stored) & \shape{(32768, 4, 64, 53)}\\
6 & answers; \code{to_channel}; skip path & \shape{(32768, 64, 32)}; \shape{(32768, 64)}; \shape{(32768, 64)}\\
7 & sum, back to a picture; BatchNorm & $\vect{z}$: \shape{(2, 64, 128, 128)}\\
\bottomrule
\end{tabular}
\caption{The reduced-attention stem of Heavy, $D=64$, token width 32, 4 heads, 40\,nm bins, at $N=185$ bands
\tagC. The weights of row 5, if stored, would be $32{,}768\times4\times64\times53=444{,}596{,}224$ numbers, 1,696\,MiB in
float32; the fused attention kernel never forms them, and 32,768 rows is exactly the stem's \texttt{SDPA\_MAX\_ROWS}, so
one call suffices.}
\label{tab:heavy:stemtrace}
\end{table}

\subsection{The encoder}

\begin{table}[htbp]
\centering\small
\begin{tabular}{@{}llll@{}}
\toprule
stage & step & shape & attention\\
\midrule
1 & patch embed $k4\,s4$, $64\to96$ & \shape{(2, 1024, 96)} & \\
 & remove 665 hidden tokens & \shape{(2, 359, 96)} & \\
 & 2 blocks; FFN $96\to384\to96$ & \shape{(2, 359, 96)} & 3 heads $\times$ 32; full, $359^2$\\
 & norm, restore, reshape & F1 \shape{(2, 96, 32, 32)} & \\\midrule
2 & embed $k3\,s2$, $96\to192$; 2 blocks & \shape{(2, 256, 192)} & 6 $\times$ 32; 16 keys\\
 & FFN $192\to768\to192$; reshape & F2 \shape{(2, 192, 16, 16)} & \\\midrule
3 & embed, $192\to480$; 2 blocks & \shape{(2, 64, 480)} & 10 $\times$ 48; 16 keys\\
 & FFN $480\to1920\to480$; reshape & F3 \shape{(2, 480, 8, 8)} & \\\midrule
4 & embed, $480\to768$; 2 blocks & \shape{(2, 16, 768)} & 12 $\times$ 64; 16 keys\\
 & FFN $768\to3072\to768$; reshape & F4 \shape{(2, 768, 4, 4)} & \\
\bottomrule
\end{tabular}
\caption{Heavy's encoder on a masked pass, $B=2$, $128\times128$, $\rho=0.65$ \tagC. ``Keys'' is the reduced key
count: $(32/8)^2=(16/4)^2=(8/2)^2=16$, and 16 tokens at stage~4.}
\label{tab:heavy:enctrace}
\end{table}

\begin{worked}[label={we:heavy:squeeze}]{What the masked stage~1 costs and holds}
Masked, stage~1 attends among the 359 kept tokens (\cref{ch:esa}): $359^2=128{,}881$ scores per head, three heads,
$3\times128{,}881=386{,}643$ per tile per block. Unmasked it would use 16 keys: $1{,}024\times16=16{,}384$ per head. The
masked pass is the expensive one, by $128{,}881/16{,}384=7.9$ times per head.

What flows where, per tile:
\begin{align*}
\text{input} &: 185\times128\times128 = 185\times16{,}384 = 3{,}031{,}040\\
\text{stem output } \vect{z} &: 64\times16{,}384 = 1{,}048{,}576\\
\text{F1} &: 96\times32\times32 = 96\times1{,}024 = 98{,}304\\
\text{F4} &: 768\times4\times4 = 12{,}288
\end{align*}
The stem cuts the values by $3{,}031{,}040/1{,}048{,}576=2.9$; stage~1 by another $1{,}048{,}576/98{,}304=10.7$, the
``10.7 to 1'' squeeze of \cref{sec:heavy:arch}.
\end{worked}

\subsection{The progressive decoder and the way back to 185 bands}

\begin{table}[htbp]
\centering\small
\begin{tabular}{@{}lp{0.45\linewidth}l@{}}
\toprule
\# & operation & shape\\
\midrule
1 & $1\times1$ projections of F1--F4 to 512; bilinear to $32\times32$ & four of \shape{(2, 512, 32, 32)}\\
2 & concatenate; fuse $2048\to512$, GELU & \shape{(2, 2048, 32, 32)} $\to$ \shape{(2, 512, 32, 32)}\\
3 & upsample $\times2$ & \shape{(2, 512, 64, 64)}\\
4 & skip: $\vect{z}$ with hidden tokens zeroed, $2\times2$ average & \shape{(2, 64, 64, 64)}\\
5 & concatenate; \texttt{up\_half}: two $3\times3$ convs, GELU & \shape{(2, 576, 64, 64)} $\to$ \shape{(2, 256, 64, 64)}\\
6 & upsample $\times2$; concatenate the full-resolution skip & \shape{(2, 256, 128, 128)}; \shape{(2, 320, 128, 128)}\\
7 & \texttt{up\_full}: two $3\times3$ convs, GELU & \shape{(2, 128, 128, 128)}\\
8 & head, $1\times1$ conv $128\to64$ & \shape{(2, 64, 128, 128)}\\
9 & per-band output weights $[\vect{u}(\lambda)\,|\,c(\lambda)]$ & \shape{(185, 65)}\\
10 & $\hat x_n=\vect{u}(\lambda_n)\cdot\hat{\vect{z}}+c(\lambda_n)$; denormalise & \shape{(2, 185, 128, 128)}\\
\bottomrule
\end{tabular}
\caption{The progressive decoder (\cref{ch:decoderwork}) and the wavelength decoder at Heavy's widths \tagC. Rows 4 and
6 are the skip that lets the decoder render detail below a token: the stem output itself, with every pixel of a
hidden token set to 0 so no hidden value reaches the reconstruction.}
\label{tab:heavy:dectrace}
\end{table}

\paragraph{The segmentation path.} The class map comes from a second, \emph{unmasked} pass through the same stem and
encoder (\code{segment}), into a \code{SegFormerDecoder} of width 256 with five outputs per sub-pixel:
four projections to \shape{(2, 256, 32, 32)}, concatenated to \shape{(2, 1024, 32, 32)}, fused to
\shape{(2, 256, 32, 32)}, refined, a sub-pixel convolution to $5\times16=80$ channels, \shape{(2, 80, 32, 32)}, and
the pixel shuffle to \shape{(2, 5, 128, 128)}: one logit per class per pixel, land being ``none of the five''
(\cref{ch:seghead}) \tagC.

\begin{figure}[htbp]\centering
\begin{tikzpicture}[node distance=3.2mm,
  box/.style={draw,rounded corners=2pt,minimum height=6mm,font=\small,align=center,text width=44mm},
  lab/.style={font=\scriptsize\ttfamily,anchor=west},
  ar/.style={-{Stealth}}]
\node[box,fill=spec!12] (in) {185 native bands};
\node[box,below=of in,fill=spec!22] (stem) {stem: 53 bins, 64 queries};
\node[box,below=of stem,fill=spec!12] (bn) {BatchNorm (64)};
\node[box,below=of bn,fill=spat!15] (e1) {stage 1: delete 665 of 1,024};
\node[box,below=of e1,fill=spat!15] (e2) {stages 2--4};
\node[box,below=of e2,fill=attn!12] (fuse) {fuse four maps to 512};
\node[box,below=of fuse,fill=attn!12] (uh) {$\times2$, + skip, \texttt{up\_half}};
\node[box,below=of uh,fill=attn!12] (uf) {$\times2$, + skip, \texttt{up\_full}, head};
\node[box,below=of uf,fill=spec!22] (out) {per-band output};
\foreach \a/\b in {in/stem,stem/bn,bn/e1,e1/e2,e2/fuse,fuse/uh,uh/uf,uf/out} \draw[ar] (\a) -- (\b);
\node[lab] at ($(in.east)+(1.5mm,0)$) {(2, 185, 128, 128)};
\node[lab] at ($(stem.east)+(1.5mm,0)$) {P = 32768 rows};
\node[lab] at ($(bn.east)+(1.5mm,0)$) {z (2, 64, 128, 128)};
\node[lab] at ($(e1.east)+(1.5mm,0)$) {(2, 359, 96) -> F1 (2, 96, 32, 32)};
\node[lab] at ($(e2.east)+(1.5mm,0)$) {F2 192/16, F3 480/8, F4 768/4};
\node[lab] at ($(fuse.east)+(1.5mm,0)$) {(2, 512, 32, 32)};
\node[lab] at ($(uh.east)+(1.5mm,0)$) {(2, 256, 64, 64)};
\node[lab] at ($(uf.east)+(1.5mm,0)$) {(2, 64, 128, 128)};
\node[lab] at ($(out.east)+(1.5mm,0)$) {(2, 185, 128, 128)};
% skip path
\draw[ar,spec,dashed] (bn.west) -- ++(-6mm,0) |- (uh.west);
\draw[ar,spec,dashed] ($(bn.west)+(-6mm,0)$) |- (uf.west);
\node[font=\scriptsize,text=spec,anchor=east,align=right] at ($(uh.west)+(-7mm,0)$) {skip:\\ hidden = 0};
% segmentation head
\node[box,fill=loss!15,text width=26mm,left=20mm of e2] (seg) {segmentation head, width 256, unmasked pass};
\draw[ar,loss] (e2.west) -- (seg.east);
\node[font=\scriptsize\ttfamily,below=1mm of seg] {(2, 5, 128, 128)};
\end{tikzpicture}
\caption{Indradhanu Heavy as served, one masked pass over two $128\times128$ PRISMA tiles with 185 bands, drawn from
the shapes the forward hooks recorded \tagC. Teal: spectral stem and output; blue: encoder; plum: progressive decoder;
dashed teal: the skip from the stem output with hidden tokens zeroed; amber: the segmentation head, fed by a separate,
unmasked encoder pass.}
\label{fig:heavy:trace}
\end{figure}

\subsection{The stem, the output and the segmentation head, counted by hand}\label{sec:heavy:stemcount}

\Cref{tab:heavy:total} took its stem rows from the built model. Count them. Every small network in the stem is a
two-layer perceptron of the 32-number wavelength encoding, $32\to64\to\text{out}$ (\file{app/foundation_models/components/spectral_reduced.py},
\code{_hyper}), so it has $(32\times64+64)+(64\times\text{out}+\text{out})=2{,}112+65\,\text{out}$ parameters.
\begin{align*}
\text{band weights } \vect{k}(\lambda),\ \text{out}=32 &: 2{,}112+65\times32 = 2{,}112+2{,}080 = 4{,}192\\
\text{skip weights } W(\lambda),\ \text{out}=64 &: 2{,}112+65\times64 = 2{,}112+4{,}160 = 6{,}272\\
\text{bin bias, } 32\to32 &: 32\times32+32 = 1{,}056\\
\text{queries} &: 64\times32 = 2{,}048\\
\text{q, k, v, o, each } 32\to32 &: 4\times1{,}056 = 4{,}224\\
\text{to\_channel, } 32\to1 &: 32+1 = 33\\
\text{stem} &: 4{,}192+6{,}272 = 10{,}464;\ +1{,}056 = 11{,}520;\ +2{,}048 = 13{,}568;\\
&\phantom{:}\ +4{,}224 = 17{,}792;\ +33 = 17{,}825\\[3pt]
\text{output weights } [\vect{u}\,|\,c],\ \text{out}=64+1 &: 2{,}112+65\times65 = 2{,}112+4{,}225 = 6{,}337\\
\text{BatchNorm on } D=64 &: 2\times64 = 128
\end{align*}
These are the 17,825, 6,337 and 128 of \cref{tab:heavy:total}; none depends on the band count. The segmentation
head is a \code{SegFormerDecoder} of width 256 writing $5\times16=80$ channels:
\begin{align*}
\text{projections} &: (96+192+480+768)\times256+4\times256 = 1{,}536\times256+1{,}024\\
&\phantom{:}\ = 393{,}216+1{,}024 = 394{,}240\\
\text{fuse} &: 1{,}024\times256+256 = 262{,}144+256 = 262{,}400\\
\text{refine } 3\times3 &: 256\times256\times9+256 = 589{,}824+256 = 590{,}080\\
\text{sub-pixel } 3\times3 &: 256\times80\times9+80 = 184{,}320+80 = 184{,}400\\
\text{head} &: 394{,}240+262{,}400 = 656{,}640;\ +590{,}080 = 1{,}246{,}720;\\
&\phantom{:}\ +184{,}400 = 1{,}431{,}120
\end{align*}
\begin{equation}\label{eq:heavy:served}
33{,}884{,}194+1{,}431{,}120 = 35{,}315{,}314
\end{equation}
trainable numbers in the served model, the documents' 35.3\,M. Loading the shipped checkpoint strictly and summing its
parameters gives 35,315,314 \tagC.

\begin{sourcenote}[: the manifest's 35,315,443]
The v2.1.0 manifest gives \code{params} 35,315,443. \file{scripts/make_model_manifest.py} counts every value in the
state dictionary except the \code{stats.*} curve and \code{trained_wl}, so it includes the stem BatchNorm's running mean
and variance ($2\times64=128$) and its batch counter (1): $35{,}315{,}314+129=35{,}315{,}443$. The difference is buffers,
not parameters, as for Indradhanu (\cref{ch:recon}) and Chakshu (\cref{ch:chakshu}).
\end{sourcenote}

\section{Pretraining on a rented H100}\label{sec:heavy:pretrain}

The run used an \emph{H100 capacity block}\index{capacity block}: AWS rents a GPU instance for a fixed window
booked in advance, here a \texttt{p5.4xlarge} (one H100 with 80\,GB, 16 vCPU) from 2026-10-28 11:30 to 2026-10-29
11:30 UTC \tagM{F002-heavy/README.md}. The window runs whether or not your job does, like a hired hall: an hour
spent setting up is an hour of the booking gone.

\begin{table}[htbp]
\centering\small
\begin{tabular}{ll}
\toprule
UTC & event\\
\midrule
10-28 11:30 & block starts\\
12:38 & launch (68 minutes into the block)\\
12:48 & first start; crashes at the first validation, three times\\
13:10 & the driver gives up after the third failure; GPU idle\\
14:05 & restart with the fix, batch $32\times1$\\
22:55 & 80 epochs done\\
10-29 01:02--10:39 & FTH1, FTB1, FTH2, FTB2, FTH3 (\cref{sec:heavy:control})\\
11:30 & block ends\\
\bottomrule
\end{tabular}
\caption{Block 1. Source: \tagM{F002-heavy/README.md} and the FTH READMEs.}
\label{tab:heavy:timeline}
\end{table}

From 11:30 to 14:05 is 2\,h\,35\,min $=155$ minutes, $155/1{,}440=0.108$ of the 24-hour booking: the README's
``\textasciitilde2.5\,h of block 1 were lost''. Split by the timeline: $68$ before launch, $10$ of start-up, $22$
in three failed starts, $55$ idle; $68+10+22+55=155$.

\subsection{The crash: 10 test shards, 12 loader workers}

Training data arrives in \emph{shards}\index{shard}: tar files of patches, read by \texttt{webdataset}. A data
\emph{loader worker}\index{loader worker} is a separate process that reads and decodes data while the GPU trains;
the configuration asked for 12. Webdataset deals shards to workers in turn, and a worker with no shard has nothing
to yield.

\begin{worked}[label={we:heavy:shards}]{Dealing shards to workers}
Deal shards $0,1,2,\dots$ round-robin: shard $i$ goes to worker $i \bmod 12$.
\begin{align*}
10\text{ shards} &: \text{workers }0\text{--}9\text{ get one each};\ \text{workers 10 and 11 get none}\\
16\text{ shards} &: \text{workers }0\text{--}3\text{ get two (shards } i,\ i+12);\ \text{workers }4\text{--}11\text{ get one}\\
\text{check} &: 4\times2 + 8\times1 = 8 + 8 = 16
\end{align*}
With 10 shards, two workers are empty and webdataset raised ``No samples found \ldots{} fewer shards than
workers''. The rule: shards $\ge$ workers. The fix was \lstinline!test_shards_per_size: 16!.
\end{worked}

Why it went unseen for about 55 minutes, in the README's words: ``My epoch watcher only looked for an epoch-1
row''. A watcher that waits for success says nothing when the job fails; it waits. The watcher was changed to
alert on any trainer exit or the driver stopping.

\subsection{The stale log line}

The second trap was in \file{h100_chain.sh}, the script that waits for pretraining and then runs the fine-tunes.
Its committed version waited for the sentence ``F002-heavy driver finished'' in a log file, then checked that
\file{metrics.csv} had 80 epochs and stopped if not. But the driver prints that sentence whenever its loop ends,
including after the three failed starts at 13:10, and the log is appended across restarts. So the \code{grep}
found the crashed run's line, counted fewer than 80 epochs, and the chain stopped itself. The fix waits on the
thing it needs, the epoch count, and stops only if the driver process is gone:
\begin{lstlisting}[language=bash]
# wait on the epoch count, not a log line: heavy.log is appended across restarts and keeps old "finished" lines
epochs() { echo $(( $(wc -l < $N/runs/F002-heavy/metrics.csv 2>/dev/null || echo 1) - 1 )); }
until [ $(epochs) -ge 80 ]; do
  pgrep -f heavy_run.sh >/dev/null || { log "driver not running at $(epochs) epochs - chain stops"; exit 1; }
  sleep 120
done
\end{lstlisting}

\begin{keypoint}
Watch for failure, not only for success, and wait on the state you need (80 rows in \texttt{metrics.csv}), never
on a message an earlier run may have left behind.
\end{keypoint}

\subsection{Batch size, throughput and the schedule}

The first start used batch 16 with 2 accumulation steps; the restart used $32\times1$. The effective batch is 32
either way. At $16\times2$ the GPU ran 63--85\% busy with 86\% of the CPUs idle; at $32\times1$, 85--100\%
\tagM{F002-heavy/README.md}. Epoch 1 took 369.4\,s, $30{,}000/369.4=81.2$ patches per second. All 80 epochs took
30,690 seconds of training, 8.5 hours.

F001 ran a 200-epoch cosine schedule (\cref{ch:pretrain}); Heavy ran 80, ``cosine fully decayed'', on the user's
``we will lose steam after epoch 50-75''. A cosine schedule (\cref{ch:trainloop}) lowers the learning rate along
half a cosine wave from $3\times10^{-4}$ to the floor $10^{-6}$:
\begin{equation}\label{eq:heavy:cosine}
\eta(e) = 10^{-6} + \left(3\times10^{-4}-10^{-6}\right)\frac{1+\cos(\pi e/80)}{2}.
\end{equation}

\begin{worked}[label={we:heavy:lr}]{The learning rate at epochs 40 and 60}
\begin{align*}
e=40 &: \cos(\pi/2) = 0;\ \tfrac{1+0}{2} = 0.5;\ 2.99\times10^{-4}\times0.5 = 1.495\times10^{-4};\ +10^{-6} = 1.505\times10^{-4}\\
e=60 &: \cos(3\pi/4) = -0.7071;\ \tfrac{1-0.7071}{2} = 0.14645;\ 2.99\times10^{-4}\times0.14645 = 4.379\times10^{-5};\\
&\phantom{:}\ +10^{-6} = 4.479\times10^{-5}
\end{align*}
\file{metrics.csv} logs $1.5050\times10^{-4}$ and $4.4788\times10^{-5}$ \tagM{F002-heavy/metrics.csv}: the schedule
is evaluated at the epoch number itself.
\end{worked}

The data rotated: training shards in blocks of 200, a new block every 5 epochs (F001: every 20), 16 blocks, 3,200
of the 3,419 training shards, 30,000 patches per epoch.

\section{The curve: did it overfit?}\label{sec:heavy:curve}

\begin{figure}[htbp]
\centering
\begin{tikzpicture}
\begin{groupplot}[group style={group size=1 by 2,vertical sep=1.3cm},width=0.9\linewidth,height=5.2cm,
  xmin=0,xmax=81,grid=major,grid style={gray!15},legend cell align=left,
  yticklabel style={/pgf/number format/fixed,/pgf/number format/precision=3}]
\nextgroupplot[ylabel={L1 (normalised)},ymin=0.015,ymax=0.042,legend pos=north east,
  scaled y ticks=false]
\addplot[loss,thick] coordinates {(1,0.04121) (2,0.03364) (3,0.03041) (4,0.02826) (5,0.02723) (6,0.02622) (7,0.02629) (8,0.02544) (9,0.02520) (10,0.02476) (11,0.02451) (12,0.02381) (13,0.02374) (14,0.02341) (15,0.02329) (16,0.02313) (17,0.02282) (18,0.02265) (19,0.02256) (20,0.02241) (21,0.02253) (22,0.02236) (23,0.02215) (24,0.02210) (25,0.02205) (26,0.02193) (27,0.02186) (28,0.02169) (29,0.02159) (30,0.02139) (31,0.02147) (32,0.02132) (33,0.02116) (34,0.02109) (35,0.02104) (36,0.02104) (37,0.02098) (38,0.02093) (39,0.02075) (40,0.02074) (41,0.02080) (42,0.02074) (43,0.02062) (44,0.02058) (45,0.02046) (46,0.02040) (47,0.02031) (48,0.02024) (49,0.02016) (50,0.02006) (51,0.01997) (52,0.01993) (53,0.01984) (54,0.01978) (55,0.01978) (56,0.01982) (57,0.01973) (58,0.01968) (59,0.01962) (60,0.01958) (61,0.01945) (62,0.01944) (63,0.01938) (64,0.01937) (65,0.01934) (66,0.01944) (67,0.01943) (68,0.01936) (69,0.01935) (70,0.01931) (71,0.01926) (72,0.01924) (73,0.01926) (74,0.01919) (75,0.01921) (76,0.01915) (77,0.01917) (78,0.01916) (79,0.01913) (80,0.01914)};
\addlegendentry{train L1}
\addplot[spec,thick] coordinates {(1,0.02854) (2,0.02709) (3,0.02403) (4,0.02286) (5,0.02221) (6,0.02064) (7,0.02069) (8,0.02100) (9,0.02004) (10,0.01960) (11,0.02029) (12,0.02085) (13,0.02052) (14,0.01964) (15,0.02078) (16,0.02088) (17,0.02036) (18,0.01976) (19,0.02005) (20,0.02016) (21,0.01942) (22,0.01877) (23,0.01868) (24,0.01873) (25,0.01839) (26,0.01900) (27,0.01837) (28,0.01856) (29,0.01807) (30,0.01784) (31,0.01898) (32,0.01858) (33,0.01914) (34,0.01777) (35,0.01797) (36,0.01801) (37,0.01767) (38,0.01856) (39,0.01774) (40,0.01761) (41,0.01719) (42,0.01729) (43,0.01730) (44,0.01740) (45,0.01767) (46,0.01709) (47,0.01754) (48,0.01749) (49,0.01711) (50,0.01681) (51,0.01736) (52,0.01745) (53,0.01705) (54,0.01694) (55,0.01729) (56,0.01742) (57,0.01722) (58,0.01703) (59,0.01706) (60,0.01714) (61,0.01680) (62,0.01668) (63,0.01669) (64,0.01668) (65,0.01657) (66,0.01657) (67,0.01640) (68,0.01646) (69,0.01637) (70,0.01672) (71,0.01631) (72,0.01619) (73,0.01629) (74,0.01627) (75,0.01624) (76,0.01644) (77,0.01631) (78,0.01635) (79,0.01634) (80,0.01636)};
\addlegendentry{val L1}
\addplot[maskgrey,dashed,domain=0:81] {0.0182};
\addlegendentry{F001 best val L1 (0.0182)}
\nextgroupplot[xlabel={epoch},ylabel={val SAM (deg)},ymin=3.5,ymax=9]
\addplot[attn,thick] coordinates {(1,8.84) (2,8.42) (3,7.37) (4,6.71) (5,6.42) (6,5.71) (7,5.57) (8,5.41) (9,5.19) (10,5.06) (11,5.09) (12,5.26) (13,5.18) (14,4.84) (15,5.30) (16,5.37) (17,5.16) (18,5.09) (19,4.97) (20,5.03) (21,4.83) (22,4.45) (23,4.58) (24,4.56) (25,4.47) (26,5.30) (27,4.56) (28,4.48) (29,4.47) (30,4.43) (31,4.97) (32,4.42) (33,4.76) (34,4.31) (35,4.51) (36,4.32) (37,4.27) (38,5.16) (39,5.16) (40,4.43) (41,4.02) (42,4.13) (43,4.19) (44,4.42) (45,4.24) (46,3.99) (47,4.32) (48,4.28) (49,4.01) (50,3.98) (51,4.28) (52,4.65) (53,4.03) (54,4.18) (55,4.13) (56,4.16) (57,4.00) (58,4.01) (59,4.06) (60,4.34) (61,4.03) (62,4.08) (63,4.01) (64,4.09) (65,3.97) (66,3.91) (67,3.88) (68,3.89) (69,3.85) (70,4.00) (71,3.84) (72,3.77) (73,3.81) (74,3.80) (75,3.78) (76,3.88) (77,3.85) (78,3.90) (79,3.92) (80,3.89)};
\addplot[maskgrey,dashed,domain=0:81] {4.06};
\addplot[only marks,mark=*,loss] coordinates {(72,3.77)};
\node[loss,anchor=south,font=\small] at (axis cs:72,3.85) {best, epoch 72};
\node[maskgrey,anchor=south west,font=\small] at (axis cs:2,4.08) {F001 best (4.06)};
\end{groupplot}
\end{tikzpicture}
\caption{Indradhanu Heavy's pretraining, all 80 epochs drawn from \texttt{metrics.csv}
\tagM{F002-heavy/metrics.csv}. Train L1 lies above validation L1 throughout (see the text), and both keep falling.}
\label{fig:heavy:curve}
\end{figure}

\begin{drift}
F002-heavy, 2026-10-28: best validation L1 / SAM 0.0162 / 3.77° at epoch 72; epoch 80 0.0164 / 3.89°. F001's best
in 113 epochs: 0.0182 / 4.06° at epoch 100 \tagM{F002-heavy/README.md}. Heavy first beat F001's L1 at epoch 29
(0.0181) and its SAM at epoch 41 (4.02°).
\end{drift}

The user's prediction holds in the numbers: from epoch 1 to 20 validation L1 fell $0.0285-0.0202=0.0083$; from 50
to 72, $0.0168-0.0162=0.0006$, fourteen times less ($0.0083/0.0006=13.8$).

One caveat the README does not make: F001 validated on test shards 0--9, Heavy on 16 test shards because of the
crash fix. Both draw 1,500 patches per validation, but not the same 1,500. The gap is large next to epoch-to-epoch
wobble, so its direction is safe; its exact size is not.

\begin{definition}[Overfitting]
A model \emph{overfits}\index{overfitting} when it learns its training examples rather than the pattern behind
them: training loss keeps falling while validation loss rises.
\end{definition}

Two pieces of evidence say Heavy did not.
\begin{enumerate}
\item \textbf{The train--validation gap is steady.} Train L1 sits above validation L1 throughout, with no trend.
  Train above validation is expected here: training hides 65\% of tokens, reads a random band subset in 80\% of
  batches and runs with dropout 0.4, while validation hides 50\% twice with complementary masks and no
  augmentation.
\item \textbf{Nothing is memorised within a block.} Each block's 200 shards are seen for 5 epochs. A model
  memorising them would show train loss falling inside a block and jumping back up at the next block. Across the
  15 block changes the largest rise is $+0.00012$ (epoch 20 to 21), and train L1 falls inside all 16 blocks;
  validation falls inside 13 of 16 \tagC.
\end{enumerate}

\begin{sourcenote}[: the size of the gap]
The course (part 68) gives the gap as 0.0023 to 0.0033 from epoch 20 to 80. Computing it from every row of
\file{F002-heavy/metrics.csv} gives 0.00201 to 0.00366 \tagC (\file{research/textbook/tools/compute_39-41_output.txt});
the course's range comes from the rounded summary rows. The conclusion, a steady gap with no trend, is unchanged.
The block-change jump ($+0.00012$) and the 16 and 13 blocks reproduce exactly.
\end{sourcenote}

\section{Pretrained detection, and the decision}\label{sec:heavy:pretraineddet}

\begin{table}[htbp]
\centering
\begin{tabular}{lccc}
\toprule
 & pAUC & caught top 1\% / 0.1\% & TPR @ 0.1\%\\
\midrule
F001 (pretrained, FT04 size) & 0.155 & 45.1\% / \textbf{14.0\%} & 0.059\\
Heavy (pretrained, epoch 72) & \textbf{0.166} & \textbf{47.4\%} / 12.4\% & 0.058\\
FT04 (F001 + fine-tuning) & 0.170 & 46.0\% / 15.8\% & 0.065\\
\bottomrule
\end{tabular}
\caption{Pretrained detection, real PRISMA bench, residual Mahalanobis, cloud excluded on FT04's masks. Source:
\tagM{F002-heavy/README.md}.}
\label{tab:heavy:pretrained}
\end{table}

Paired bootstrap, Heavy $-$ F001: pAUC $+0.0111\ [-0.0003,+0.0223]$, $P=0.971$: the interval touches zero, not
resolved, by a hair. Caught at the top 0.1\%: $-1.5$ points $[-3.0,-0.2]$, a resolved loss.

\begin{worked}[label={we:heavy:set4}]{Where the top-0.1\% loss sits}
Per-scene pAUC, Heavy 0.259 / 0.150 / 0.071 / 0.185, F001 0.220 / 0.126 / 0.080 / 0.195:
\begin{align*}
\text{Heavy} &: 0.259+0.150+0.071+0.185 = 0.665;\quad 0.665/4 = 0.166\\
\text{F001} &: 0.220+0.126+0.080+0.195 = 0.621;\quad 0.621/4 = 0.155
\end{align*}
Objects caught at the top 0.1\%, Heavy 53 / 23 / 12 / 25, F001 52 / 27 / 11 / 37:
\begin{align*}
\text{Heavy} &: 53+23+12+25 = 113;\quad 113/909 = 0.124\\
\text{F001} &: 52+27+11+37 = 127;\quad 127/909 = 0.140\\
\text{loss} &: 127-113 = 14 \text{ objects, of which set4 } 37-25 = 12
\end{align*}
Twelve of the fourteen lost objects are on set4, the hazy scene. FT04's lead over F001 is
$0.170-0.155=0.015$; pretrained Heavy recovers $0.166-0.155=0.011$ of it, $0.011/0.015=0.73$ (the README rounds
to ``\textasciitilde75\%'').
\end{worked}

The decision, 2026-10-29: ``Indradhanu Heavy is the base pretrained model'' (user: ``I think heavy is our base
model now''). Be clear what the evidence was at that moment: better reconstruction, a near-resolved pAUC gain as a
pretrained model, and a resolved loss at the top 0.1\%. It was a judgement on the better general representation,
not a detection win. That came next.

\section{The fair test: a control and three checkpoints}\label{sec:heavy:control}

FT04 is F001 plus a four-stage curriculum (\cref{ch:finetune,ch:ftruns}). Block 1 had about 12 hours left after
pretraining, not enough to repeat four stages for Heavy. So Heavy got FT04's \emph{final} recipe in one joint run.
If that model beat FT04, you could not say whether the bigger model won or something else in the change did; if it
lost, it might be the short recipe. Two causes mixed into one comparison: that is
\emph{confounding}\index{confounding}.

\begin{definition}[Control]
A \emph{control}\index{control (experiment)} is a run identical to the experiment in every way except the thing
under test.
\end{definition}

\file{FTB1-base-joint/config.json} says it directly: ``F001 (FT04's size, SegFormer decoder) fine-tuned with exactly
FTH1's one-run recipe, so Heavy vs FTB1 isolates size + decoder + pretraining from the recipe.'' The control
separates the \emph{model} from the \emph{recipe}; it does not separate size alone, because Heavy also differs in
decoder, pretraining run and segmentation-head width (256 against 128).

\begin{center}\small
\begin{tabular}{lllll}
\toprule
run & starts from & joint epochs & lr (cosine) & total joint epochs\\
\midrule
FTH1 & Heavy, epoch 72 & 10 & $10^{-4}$ & 10\\
FTB1 (control) & F001 epoch 109 & 10 & $10^{-4}$ & 10\\
FTH2 & FTH1 & 8 & $5\times10^{-5}$ & 18\\
FTB2 (control) & FTB1 & 8 & $5\times10^{-5}$ & 18\\
FTH3 & FTH2 & 6 & $3\times10^{-5}$ & 24 (no control)\\
\bottomrule
\end{tabular}
\end{center}

\Cref{ch:decoderwork} ended on a lesson: FT09b's gain cleared the object bootstrap and vanished five epochs later.
So these runs kept a checkpoint after every epoch, and every comparison is made at \textbf{three neighbouring
checkpoints}. A gain counts if it holds at all three.

\subsection{FTH1 and FTB1}

At epoch 10 both one-run models fall short of FT04's curriculum on PRISMA reconstruction and masks, and Heavy
reconstructs \emph{worse} than the control: SAM 3.94 / 4.74 / 5.60° against FTB1's 3.90 / 4.48 / 5.34° (FT04 3.88 /
4.11 / 5.46°) \tagM{FTH1-heavy-joint/README.md}. On these numbers alone Heavy looked like a disappointment.

Detection, with cloud excluded on FT04's masks for both, so only the detector differs:
\begin{center}\small
\begin{tabular}{lccc}
\toprule
checkpoint & FTB1 pAUC & FTH1 pAUC & FTH1 $-$ FTB1 [95\%]\\
\midrule
epoch 8 & 0.162 & \textbf{0.175} & $+0.0132\ [+0.0022,+0.0243]$\\
epoch 9 & 0.162 & \textbf{0.171} & $+0.0093\ [-0.0013,+0.0195]$\\
epoch 10 & 0.161 & \textbf{0.175} & $+0.0145\ [+0.0045,+0.0246]$\\
\bottomrule
\end{tabular}
\end{center}
Heavy ahead at all three, resolved at two. \textbf{Worse reconstruction, better detection.}

\begin{sourcenote}[: 0.182 or 0.183]
Against shipped FT04 with the control's masks, the FTH1 README gives epochs 8 / 9 / 10 pAUC 0.182 / 0.178 / 0.182.
Its \file{eval_report.md} row for epoch 8 reads 0.183 (\file{FTH1-heavy-joint/eval_report.md}, line 17). The two
are the same number rounded differently; the bootstrap's $+0.0130$ over FT04's 0.1695 puts it at 0.1825.
\end{sourcenote}

\section{The mask failure: haze in the thin-cloud head}\label{sec:heavy:haze}

Excluding cloud with FTH1's \emph{own} head wrecked every detector it was applied to: FT04 fell from 0.170 to 0.127
pAUC with FTH1's masks, and FTH1 itself from 0.182 to 0.133 (ROC-AUC 0.742, caught top 1\% 32.8\%)
\tagM{FTH1-heavy-joint/eval_report.md}. Yet FTH1's EnMAP validation cloud IoU, 0.717, was close to FTB1's 0.721.

The product's class map comes from \code{segment_full_scene}. The head predicts five independent probabilities
per pixel (water, cloud, thin cloud, shadow, snow; \cref{ch:seghead}), and the class map sets priorities by
assignment order, cloud last so it wins:
\begin{lstlisting}
cls[prob[0] > 0.5] = 1                              # water
cls[prob[3] > 0.5] = 3                              # shadow
cls[prob[4] > 0.5] = 4                              # snow
cls[np.maximum(prob[1], prob[2]) > 0.5] = 2         # cloud: cloud OR thin cloud
\end{lstlisting}
A pixel is excluded as cloud if \emph{either} cloud output passes 0.5.

\begin{table}[htbp]
\centering
\begin{tabular}{lcccc}
\toprule
 & set1 & set2 & set3 & set4\\
\midrule
FT04 & 0.0\% / 0 & 0.0\% / 0 & 0.8\% / 0 & 11.8\% / 2\\
FTB1 & 0.0\% / 0 & 0.0\% / 0 & 1.7\% / 0 & 0.3\% / 0\\
FTH1 & 0.4\% / 17 & 0.0\% / 0 & 2.1\% / 0 & \textbf{75.4\% / 209 of 230}\\
\bottomrule
\end{tabular}
\caption{Share of each bench scene called cloud, and targets inside it. Source: \tagM{FTH1-heavy-joint/README.md}.}
\label{tab:heavy:cloudshare}
\end{table}

``All of it is the thin-cloud output'': on set4, FTH1's thick-cloud channel covers 0.4\% (no targets), its
thin-cloud channel 75.4\% (209 targets). That is the haze over set4. Thin cloud was always the weak label (its IoU
was 0.07 for FTH1), and the validation IoU above measures the \emph{cloud} channel, which is why it did not warn
anyone.

\subsection{Why zeroing targets wrecks pAUC}

Cloud exclusion sets the score of every excluded pixel to 0 after scoring. Recall from
\cref{ch:measures,ch:bootstrap} the per-target credit $c(p)=\max(0,\ 0.01-\mathrm{FPR}(p))/0.01$, where
$\mathrm{FPR}(p)$ is the share of background pixels scoring at or above target pixel $p$, and pAUC is the mean
credit. A target scored 0 has every background pixel at or above it: $\mathrm{FPR}=1$, so
$c=\max(0,0.01-1)/0.01=0$.

\begin{worked}[label={we:heavy:zero}]{Switching off a scene's targets}
Ten one-pixel targets, each with $c=0.5$. Zero nine of them.
\begin{align*}
\text{before} &: (10\times0.5)/10 = 5.0/10 = 0.50\\
\text{after} &: (1\times0.5 + 9\times0)/10 = 0.5/10 = 0.05
\end{align*}
Now set4, as an estimate. FTH1 epoch 10 with thick-cloud-only exclusion (below) scores set4 at 0.222. Assume
(\tagA) the 209 excluded targets are typical of set4's 230, so the 21 kept ones carry their share:
\begin{align*}
\text{kept share} &: 21/230 = 0.0913\\
\text{set4 pAUC} &: 0.222\times0.0913 = 0.020\\
\text{bench, thick} &: (0.253+0.162+0.086+0.222)/4 = 0.723/4 = 0.181\\
\text{bench, estimate} &: (0.253+0.162+0.086+0.020)/4 = 0.521/4 = 0.130\\
\text{measured} &: 0.133
\end{align*}
The estimate ignores set1's 17 excluded targets and assumes the targets alike, so do not read the third decimal.
But it shows the collapse is one scene's targets being switched off.
\end{worked}

\begin{keypoint}
A mask head that calls haze ``cloud'' deletes the targets under the haze.
\end{keypoint}

\subsection{The fix: thick cloud only}

The evaluator gained a switch, \code{ModelSpec.cloud_rule} in \file{app/evaluation/config.py}:
\code{"cloud_or_thin"} (the default) or \code{"thick"}, described as ``thin cloud fires on haze''. Under the thick
rule a pixel that was cloud only because of thin cloud falls to whatever it would otherwise have been: snow,
shadow, water or land. With FTH1's own head and the thick rule: epoch 10 pAUC 0.181 (sets 0.253 / 0.162 / 0.086 /
0.222), the same as with FTB1's masks. Does the thick rule hurt FT04? FT04 shipped $\to$ FT04 thick:
$-0.0004\ [-0.0015,+0.0015]$, a tie \tagM{FTH2-heavy-joint-more/bootstrap.txt}.

\begin{sourcenote}[: where the thick rule lives now]
The course (part 69) shows the rule as an \code{np.select} in \file{app/evaluation/models.py} that re-assigns
thin-only pixels after \code{segment_full_scene}. In the current source the rule has moved into
\code{segment_full_scene} itself, as a parameter \code{thin_cloud} (default \code{True}): the cloud line becomes
\lstinline!cls[(np.maximum(prob[1], prob[2]) if thin_cloud else prob[1]) > 0.5] = 2!, and \file{models.py} calls it with
\code{thin_cloud=self.spec.cloud_rule != "thick"}
(\file{app/foundation_models/inferencers/spectral_stem_inferencer.py}, lines 112--141). The two give the same class
map: with thin cloud ignored, a thin-only pixel keeps the class from the earlier lines, where snow overwrites
shadow overwrites water, which is the course's priority snow, shadow, water, land.
\end{sourcenote}

\section{FTH2, FTB2 and FTH3: a gain that holds}\label{sec:heavy:fth2}

From here every model excludes cloud with its own thick-cloud masks.

\begin{table}[htbp]
\centering\small
\begin{tabular}{lccc}
\toprule
 & pAUC epoch 6 / 7 / 8 & caught top 1\% & caught top 0.1\%\\
\midrule
FT04 shipped (cloud or thin) & 0.170 & 46.0\% & \textbf{15.8\%}\\
FTB2 (own thick) & 0.160 / 0.163 / 0.162 & 46.1 / 46.4 / 46.8\% & 15.1 / 15.1 / 15.1\%\\
FTH2 (own thick) & \textbf{0.185 / 0.184 / 0.186} & \textbf{48.2 / 49.0 / 49.0\%} & 15.3 / 15.2 / 15.7\%\\
\bottomrule
\end{tabular}
\caption{18 joint epochs. Source: \tagM{FTH2-heavy-joint-more/README.md}.}
\label{tab:heavy:fth2}
\end{table}

Every pAUC interval clears zero, at every checkpoint, against both the control and FT04
(\cref{fig:heavy:intervals}). The average gap over the control:
\begin{align*}
(0.0245+0.0209+0.0244)/3 = 0.0698/3 = 0.0233 .
\end{align*}

\begin{sourcenote}[: ``about twice'' or ``about three times'']
The course (part 69) quotes the FTH2 README as calling this gap ``about twice what cloud exclusion bought'', and
corrects it: cloud exclusion is $+0.0074$ (\cref{ch:bootstrap}), and $0.0233/0.0074=3.1$. The current README
reads ``about three times what cloud exclusion bought, +0.0074''
(\file{FTH2-heavy-joint-more/README.md}): the source was corrected after the course note.
\end{sourcenote}

The gap \emph{grew} with fine-tuning: about $+0.013$ at 10 epochs, $+0.023$ at 18. And reconstruction still did not
separate the models: FTH2's PRISMA SAM (4.48°) is worse than the control's (4.36°).

\paragraph{FTH3: the plateau.} Six more joint epochs at $3\times10^{-5}$, Heavy only, from the source
\tagM{FTH3-heavy-joint-more2/README.md}. Epochs 4, 5 and 6 score pAUC 0.184, 0.189 and 0.187. Against FTH2 epoch 8
the three differences are:
\begin{center}
$-0.0021\ [-0.0043,-0.0000]$,\qquad $+0.0028\ [-0.0002,+0.0058]$,\qquad $+0.0011\ [-0.0009,+0.0030]$:
\end{center}
one barely below zero, one barely above, one a tie. That is a \emph{plateau}: more epochs at this recipe no longer move detection. Against
FT04 all three stay resolved: $+0.0148$, $+0.0196$, $+0.0179$. The release candidate is FTH3 epoch 5: pAUC 0.189
against 0.170, caught top 1\% 49.4\% against 46.0\%, top 0.1\% 16.0\% against 15.8\% (a tie, $[-1.2,+1.5]$). Pick it
as the plateau's representative, not as a discovery: it is the best of a stable set and the only one whose
top-0.1\% count is not below FT04's.

\begin{worked}[label={we:heavy:mean}]{Mean gain over FT04 at three checkpoints}
FTH2 $-$ FT04 at epochs 6, 7, 8: $+0.0151$, $+0.0143$, $+0.0168$.
\begin{align*}
0.0151+0.0143+0.0168 &= 0.0462\\
0.0462/3 &= 0.0154
\end{align*}
The lowest lower bound is $+0.0032$ (epoch 7), so none touches zero. FTH3 at epochs 4, 5, 6:
$0.0148+0.0196+0.0179=0.0523$, $0.0523/3=0.0174$.
\end{worked}

\begin{figure}[htbp]
\centering
\begin{tikzpicture}
\begin{axis}[width=0.9\linewidth,height=9cm,xmin=-0.004,xmax=0.04,ymin=0.3,ymax=13.7,
  ytick={1,...,13},yticklabels={FTH3 ep6 $-$ FT04,FTH3 ep5 $-$ FT04,FTH3 ep4 $-$ FT04,FTH2 ep8 $-$ FT04,FTH2 ep7 $-$ FT04,FTH2 ep6 $-$ FT04,FTH2 ep8 $-$ FTB2,FTH2 ep7 $-$ FTB2,FTH2 ep6 $-$ FTB2,FTH1 ep10 $-$ FTB1,FTH1 ep9 $-$ FTB1,FTH1 ep8 $-$ FTB1,cloud excl. (FT04)},
  yticklabel style={font=\scriptsize},xlabel={pAUC difference [95\% paired-bootstrap interval]},
  xtick={0,0.01,0.02,0.03,0.04},
  scaled x ticks=false,xticklabel style={/pgf/number format/fixed,/pgf/number format/precision=2,font=\scriptsize},
  grid=major,grid style={gray!15}]
\draw[thick,maskgrey] (axis cs:0,0.3) -- (axis cs:0,13.7);
\addplot[only marks,mark=*,loss,error bars/.cd,x dir=both,x explicit] coordinates {
 (0.0179,1) += (0.0113,0) -= (0.0124,0)
 (0.0196,2) += (0.0120,0) -= (0.0124,0)
 (0.0148,3) += (0.0111,0) -= (0.0119,0)
 (0.0168,4) += (0.0107,0) -= (0.0115,0)
 (0.0143,5) += (0.0104,0) -= (0.0111,0)
 (0.0151,6) += (0.0106,0) -= (0.0115,0)};
\addplot[only marks,mark=square*,spat,error bars/.cd,x dir=both,x explicit] coordinates {
 (0.0244,7) += (0.0118,0) -= (0.0109,0)
 (0.0209,8) += (0.0112,0) -= (0.0107,0)
 (0.0245,9) += (0.0107,0) -= (0.0105,0)
 (0.0145,10) += (0.0101,0) -= (0.0100,0)
 (0.0093,11) += (0.0102,0) -= (0.0106,0)
 (0.0132,12) += (0.0111,0) -= (0.0110,0)};
\addplot[only marks,mark=diamond*,spec,error bars/.cd,x dir=both,x explicit] coordinates {
 (0.0074,13) += (0.0027,0) -= (0.0027,0)};
\end{axis}
\end{tikzpicture}
\caption{Heavy against the control (indigo squares) and against shipped FT04 (amber circles), at three
neighbouring checkpoints each, drawn to scale; cloud exclusion (teal) for reference. Only FTH1 epoch 9 touches zero.
Sources: \tagM{FTH1-heavy-joint/README.md}, \tagM{FTH2-heavy-joint-more/bootstrap.txt},
\tagM{FTH3-heavy-joint-more2/bootstrap.txt}.}
\label{fig:heavy:intervals}
\end{figure}

\section{The lesson: reconstruction did not separate the models}\label{sec:heavy:lesson}

\begin{center}
\begin{tabular}{lcc}
\toprule
 & PRISMA SAM (deg) & pAUC\\
\midrule
FTB1 (control, 10 epochs) & \textbf{4.48} & 0.161--0.162 (FT04 masks)\\
FTH1 (Heavy, 10 epochs) & 4.74 & \textbf{0.171--0.175} (FT04 masks)\\
FTB2 (control, 18 epochs) & \textbf{4.36} & 0.160--0.163 (own thick masks)\\
FTH2 (Heavy, 18 epochs) & 4.48 & \textbf{0.184--0.186} (own thick masks)\\
\bottomrule
\end{tabular}
\end{center}

\textbf{The model that reconstructs worse is the better detector, at every stage.} Had the decision been made on
validation SAM, the number every run logs each epoch, Heavy would have been dropped after FTH1. The FTH2 README:
``Better detection without better reconstruction.'' This is the third time in Part VIII: FT06's best-NLL epoch was
its worst for detection, FT09c reconstructed better and detected worse, and now this.

\begin{keypoint}
Judge detectors on detection: reconstruction error and detection come apart, so a model must be selected on the
measure the product uses.
\end{keypoint}

\section{Into the product}\label{sec:heavy:product}

FTH3's verdict listed three things the release candidate needed before an analyst could use it: the thick-cloud
rule, a manifest and resolver path for the progressive decoder, and a check of the worker's CPU cost. All three were
done on 2026-10-29, and FTH3 epoch 5 became Indradhanu~V2 \texttt{v2.1.0} the same day, behind FT04's codename and
architecture key, with FT04 (\texttt{v2.0.0}) kept as the manifest's first alternative
(\file{docs/lld/spectral-stem.md}, ``Serving''). Because the manifest embeds the checkpoint's whole model
configuration (\cref{ch:shipping}), switching models changed files, not code.

\paragraph{What the product loads.} The serving documents give 35.3\,M parameters. \Cref{sec:heavy:width} counted the
pretrained network at 33,884,194; the fine-tuned checkpoint adds the segmentation head, 1,431,120 counted by hand in
\cref{sec:heavy:stemcount}, for the 35,315,314 of \cref{eq:heavy:served}: the 35.3\,M \tagC.

\paragraph{The thick-cloud rule in the product.} The resolver's capabilities entry for \code{spectral_stem_finetune}
sets \code{thin_cloud_is_cloud=False}; the worker passes it to \lstinline!segment_full_scene(..., thin_cloud=False)!, so
only the head's thick-cloud channel can make a pixel ``cloud'', and \texttt{stats.json} records
\texttt{cloud\_rule: "thick"} (\cref{ch:shipping}, the source check on the cloud rule). Pixels that are only thin
cloud fall to the next class by priority and are scored.

\begin{worked}[label={we:heavy:parity}]{Parity, to the digit}
The parity run (\file{FTH3-heavy-joint-more2/product_parity.py}) takes the product's own path, resolver manifest to
inferencer to \code{read_bands}, prediction, residual Mahalanobis, segmentation with the thick rule and the zeroing,
and compares with the evaluator's run on the four PRISMA scenes \tagM{research/experiments/indradhanu-finetune/FTH3-heavy-joint-more2/README.md}:
\begin{center}\small
\begin{tabular}{lcccc}
\toprule
scene & thick cloud excluded & pAUC product & pAUC evaluator & gap\\
\midrule
set1 & 0.00\% & 0.2617 & 0.26165 & $+0.00005$\\
set2 & 0.00\% & 0.1874 & 0.18742 & $-0.00002$\\
set3 & 1.49\% & 0.0861 & 0.08609 & $+0.00001$\\
set4 & 0.47\% & 0.2217 & 0.22167 & $+0.00003$\\
\bottomrule
\end{tabular}
\end{center}
The product column's mean:
\begin{align*}
0.2617+0.1874 &= 0.4491\\
0.4491+0.0861 &= 0.5352\\
0.5352+0.2217 &= 0.7569\\
0.7569/4 &= 0.1892 .
\end{align*}
The evaluator's: $0.26165+0.18742+0.08609+0.22167=0.75683$, and $0.75683/4=0.1892$. Every gap is smaller than
half of the product column's last printed digit, $0.00005$: the gaps are rounding. Note set4: 0.47\% thick cloud,
where the old rule had called 75.4\% of the scene cloud (\cref{sec:heavy:haze}).
\end{worked}

In objects: the comparison page shown to the user on 2026-10-29 (\url{https://claude.ai/artifact/2dZFbGN6HhMBCx3kkZ8pQ4})
counts the 909 labelled targets of the four scenes caught in each model's top 1\%, FT04 418 and Heavy 449.
\begin{align*}
418/909 &= 0.4598 \to 46.0\%, & 449/909 &= 0.4939 \to 49.4\%, & 449-418 &= 31 \text{ objects,}
\end{align*}
the FTH3 README's caught-at-1\% figures for FT04 and FTH3 epoch 5 \tagC.

\paragraph{Sixteen threads.} The first CPU check (\file{cpu_timing.py}, inside the worker container, PRISMA set1,
$1186\times1196$ pixels, 185 bands) gave Heavy 1,819\,s and FT04 1,058\,s, about 30 and 18 minutes a scene. The
worker was running torch on 2 threads. A \emph{thread}\index{thread} is one line of work a program runs at a time;
a numerical library splits a large matrix product across as many threads as it is allowed. The cap of 2 was set in
two places, the worker image's \code{ENV} (\file{backend/Dockerfile.worker}) and an
\lstinline!os.environ.setdefault(..., "2")! in \file{backend/allotrope_worker/__init__.py}, after crashes that the
compose file's comment says were ``Docker Desktop UI's restart button, not throttling''. Since 2026-10-29
\file{docker/docker-compose.yml} sets \code{OMP_NUM_THREADS}, \code{MKL_NUM_THREADS}, \code{OPENBLAS_NUM_THREADS} and
\code{NUMEXPR_NUM_THREADS} to \code{${WORKER_THREADS:-16}}: compose's \code{environment} replaces the image's
\code{ENV}, and \code{setdefault} never overwrites a variable that is already set, so the 16 wins inside compose and
the 2 still applies outside it.

\begin{worked}[label={we:heavy:threads}]{Four times faster from eight times the threads}
From the FTH3 README \tagM{research/experiments/indradhanu-finetune/FTH3-heavy-joint-more2/README.md}:
\begin{center}\small
\begin{tabular}{lrrr}
\toprule
step & 2 threads & 16 threads & speed-up\\
\midrule
Heavy, total & 1{,}819\,s & 454\,s & $1819/454=4.01$\\
\quad predict & 1{,}629\,s & 402\,s & $1629/402=4.05$\\
\quad segment & 185\,s & 47\,s & $185/47=3.94$\\
\quad score & 4\,s & 4\,s & $1.00$\\
FT04, total & 1{,}058\,s & 274\,s & $1058/274=3.86$\\
\bottomrule
\end{tabular}
\end{center}
In minutes, Heavy goes from $1819/60=30.3$ to $454/60=7.6$, and still costs $454/274=1.66$ times FT04.
Eight times the threads ($16/2=8$) bought four times the speed. A standard model for why is \emph{Amdahl's
law}\index{Amdahl's law}: if a share $p$ of the work splits perfectly over $n$ threads and the rest does not, the time
is proportional to $(1-p)+p/n$. It is a model, not a measurement. Set the speed-up from 2 to 16 threads to 4 and
solve for $p$:
\begin{align*}
\frac{(1-p)+p/2}{(1-p)+p/16} &= 4 &&\text{the measured ratio}\\
1-\tfrac12p &= 4\big(1-\tfrac{15}{16}p\big) &&\text{simplify each side, multiply up}\\
1-0.5p &= 4-3.75p\\
3.25p &= 3 &&\text{collect the $p$ terms}\\
p &= 0.923 .
\end{align*}
The measured speed-up behaves as if about 92\% of the work were perfectly parallel \tagC. The score step did not speed up
at all, but at 4\,s it is a small part; the rest of the serial share is reading, writing and bookkeeping, and threads also
share memory bandwidth, which the law ignores.
\end{worked}

\begin{sourcenote}[: Heavy has shipped]
The course's part 69, as written on 2026-10-29, ended with ``Not shipped: the product still runs Indradhanu V2 (FT04)''
(corrected on 2026-10-31; the course's part 70 now covers the ship). The source had moved on the same day:
FTH3 epoch 5 was shipped on 2026-10-29 as Indradhanu V2 \texttt{v2.1.0}, same codename and architecture key as FT04
(FT04 \texttt{v2.0.0} is the manifest's first alternative), and the product now excludes thick cloud only (resolver
\code{thin_cloud_is_cloud=False}) (\file{FTH3-heavy-joint-more2/README.md}, ``Shipped''; commit c4b3a63). The three items FTH3's verdict listed before
release (the thick-cloud rule, a manifest path for the progressive decoder, a worker CPU-cost check) were done.
\end{sourcenote}

\begin{drift}
Product checks recorded 2026-10-29 in \file{FTH3-heavy-joint-more2/README.md}. Parity: the product path scored
PRISMA sets 1--4 at pAUC 0.2617 / 0.1874 / 0.0861 / 0.2217 against the evaluator's 0.26165 / 0.18742 / 0.08609 /
0.22167, mean 0.1892 for both. Worker CPU cost on PRISMA set1 ($1186\times1196\times185$ bands, 2 torch threads):
Heavy 1,819\,s against FT04 1,058\,s, $1{,}819/1{,}058=1.7$ times, far less than the 6.1 times parameter ratio;
on 16 threads (compose, since 2026-10-29) 454\,s against 274\,s \tagM{FTH3-heavy-joint-more2/README.md}. The
Dockerfile and \file{allotrope_worker/__init__.py} still default to 2 for anything run outside compose. The open caveats remain: four scenes, and a bootstrap that resamples
objects, not scenes.
\end{drift}

\begin{confusions}
\confusion{Heavy is 4x wide.}{That was its launch name. Its encoder widths are 3 times FT04's, its latent and
decoder 2 times; its parameter count is 6.1 times overall and 8.9 times for the encoder.}
\confusion{Three times the width should mean three times the parameters.}{Almost every weight is in a $C\times C$
matrix, so the count grows like $C^2$: $3^2=9$, less a little for the biases and norms (8.86 at stage 1).}
\confusion{Six times the parameters should mean six times slower.}{On the H100 an epoch took about 370\,s; on the
product's CPU worker Heavy took 1.7 times FT04's time, on 2 threads and on 16 alike. Speed depends on hardware and on how busy it is kept, not
on the parameter count alone.}
\confusion{FTH1 $-$ FTB1 measures the effect of size.}{It measures the effect of the model: size, decoder,
pretraining run and head width together. The recipe is what it holds fixed.}
\confusion{FTH1's masks were bad, so Heavy's segmentation head is broken.}{Its thick-cloud channel works;
thick-only exclusion with its own head matches the control's masks. The thin-cloud channel fires on haze, and the
product's old rule took either channel. FT04's head did it too, on 11.8\% of set4.}
\confusion{The thick rule was chosen to make Heavy look good.}{It was chosen after a diagnosis showed thin cloud
covering 75\% of a hazy scene, and checked on FT04, where it changes nothing ($-0.0004\ [-0.0015,+0.0015]$). It was
still adopted after seeing results, one more reason the release needs more than this bench.}
\confusion{FTH3 epoch 5 beat FTH2, so the extra epochs helped.}{Its gain over FTH2 is unresolved and its neighbour
is slightly below: a plateau.}
\confusion{The served model drops 20 of PRISMA's 185 bands.}{A pretrained model does; its record of trained
wavelengths is the SpectralEarth curve. The served v2.1.0's record grew to 733 wavelengths during fine-tuning, and it
keeps all 185.}
\confusion{35,315,443 parameters.}{That is the manifest's count, which includes 129 BatchNorm buffers; the trainable
count is 35,315,314.}
\confusion{Train loss above validation loss means something is broken.}{Training is the harder task here: more
hidden tokens, missing bands, dropout. Watch the gap's trend, not its sign.}
\end{confusions}

\begin{checkyourself}
\item\label{ex:heavy:1} Use \cref{eq:heavy:block} to count FT04's and Heavy's stage-2 block ($C=64$ and $192$,
$R=4$). Check both against \cref{tab:heavy:blocks}, and compute the ratio. Then predict the ratio for a 2 times
wider FT04 at stage 1.
\item\label{ex:heavy:2} Count Heavy's stage-2 patch embedding ($3\times3$, stride 2, 96 to 192 channels, plus
LayerNorm), and FT04's ($32\to64$).
\item\label{ex:heavy:3} Using \cref{eq:heavy:cosine}, what learning rate does epoch 20 log? ($\cos(\pi/4)=0.7071$.)
\item\label{ex:heavy:4} With 16 test shards, what is the largest number of loader workers that avoids an empty
worker? With 12 workers, how many shards would give every worker exactly two?
\item\label{ex:heavy:5} A pixel's thick-cloud probability is 0.3, thin cloud 0.8, water 0.6, and nothing else passes
0.5. What class does the old rule give it, and what does the thick rule give it?
\item\label{ex:heavy:6} A scene has 8 one-pixel targets with credits 1.0, 0.8, 0.8, 0.6, 0.4, 0.4, 0.0, 0.0. Compute
its pAUC; then exclusion zeroes the first three. Compute it again. Why is the drop larger than 3 of 8?
\item\label{ex:heavy:7} Why could FTH1 losing to FT04 not have shown that Heavy is the worse model, and what does
FTB1 fix and not fix?
\item\label{ex:heavy:8} With $p=0.923$, predict the speed-up from 2 to 24 threads with Amdahl's law. Then give two
reasons not to trust the prediction on the 24-core host.
\item\label{ex:heavy:9} Count by hand the stem's output network (\code{spectral_out}) for FT04's latent width $D=32$,
and check it against \cref{tab:heavy:total}. Then count FT04's segmentation head: width 128, five classes, FT04's
encoder widths 32, 64, 160, 256.
\item\label{ex:heavy:10} At inference the product hides half the stage-1 tokens of each tile, twice. For Heavy on a
$128\times128$ tile, how many tokens does each pass keep, and how many attention scores per head does stage~1 compute?
With the product's batch of 8 tiles, how many rows does the stem see, and how many calls of the fused attention
kernel does that take with \texttt{SDPA\_MAX\_ROWS} $=32{,}768$?
\end{checkyourself}

\par\bigskip\noindent\textit{That closes the experiments. Indradhanu Heavy is the base model, and its fine-tuned
FTH3 epoch 5 is what the product's Indradhanu V2 now runs; every number behind that decision is in this chapter,
with its interval.}

\nextchapter{ch:namer}
